\documentclass[11pt]{article}
\usepackage[T1]{fontenc}
\usepackage{amsmath,amssymb,amsthm,mathtools,geometry,hyperref,array,xcolor}
\geometry{margin=1in}
\setlength{\emergencystretch}{3em}
\hypersetup{hidelinks}
\newtheorem{theorem}{Theorem}[section]
\newtheorem{lemma}[theorem]{Lemma}
\newtheorem{proposition}[theorem]{Proposition}
\theoremstyle{remark}\newtheorem{remark}[theorem]{Remark}
\newcommand{\R}{\mathbb R}
\newcommand{\M}{\mathcal M}
\newcommand{\curl}{\operatorname{curl}}
\newcommand{\Div}{\operatorname{div}}
\newcommand{\supp}{\operatorname{supp}}
\newcommand{\norm}[2]{\left\|#1\right\|_{#2}}
\newcommand{\dd}{\,dx}
\title{Global weak Euler flows for finite signed vorticity on the plane}
\author{Research manuscript}
\date{8 October 2026}
\begin{document}
\maketitle

\noindent Research proof manuscript.
Verification status is recorded in the accompanying audit report;
this source alone is not a verification certificate.

\paragraph{Exact target.}
Let $K(x)=(-x_2,x_1)/(2\pi|x|^2)$. Given any finite signed Radon measure
$\mu_0\in\M(\R^2)\cap H^{-1}(\R^2)$, with the \emph{inhomogeneous} Sobolev
space and no sign restriction, set $u_0=K*\mu_0$. The requested conclusion is a
global weak Euler solution
\[
 u\in C_w([0,\infty);L^2_{\rm loc}(\R^2;\R^2)),\qquad u(0)=u_0,
 \qquad \Div u=0,
\]
such that, at \emph{every} $t\geq0$,
\[
 \mu(t)=\curl u(t)\in\M(\R^2),\qquad u(t)=K*\mu(t),\qquad
 \sup_{0\leq t\leq T}\norm{\mu(t)}{\M}<\infty\quad(T<\infty).
\]
The velocity identity is precisely
\[
 \int_0^T\!\!\int_{\R^2}
 \bigl(u\cdot\partial_t\varphi+(u\otimes u):\nabla\varphi\bigr)\dd\,dt
 +\int_{\R^2}u_0\cdot\varphi(0)\dd=0
\]
for all divergence-free $\varphi\in C_c^\infty([0,T)\times\R^2;\R^2)$.
No global kinetic-energy assumption, compact support, zero circulation, or
energy inequality is added.

\section{Initial normalization and a bounded Euler background}

We use $\nabla^\perp=(-\partial_2,\partial_1)$,
$(\nabla v)_{ij}=\partial_jv_i$, $(\Div S)_i=\sum_j\partial_jS_{ij}$, and
$\widehat f(\xi)=\int e^{-ix\cdot\xi}f(x)\dd$.
The multiplier of $K$ is $m(\xi)=-i\xi^\perp/|\xi|^2$.

\begin{lemma}[Canonical velocity and the physical-space decomposition]\label{lem:decomp}
Let $\mu\in\M\cap H^{-1}$ and $M=\norm\mu\M$. The convolution $K*\mu$
is defined almost everywhere, belongs to $L^2_{\rm loc}$, and has divergence
zero and curl $\mu$. Fix a nonnegative radial $a\in C_c^\infty(B_A)$ of
integral one. Set
\[
 \nu_0=a*\mu,\quad \sigma_0=\mu-\nu_0,\quad
 V_0=K*\nu_0,\quad v_0=(K-K*a)*\mu.
\]
Then $K*\mu=V_0+v_0$, $v_0\in L^1\cap L^2$, $V_0\in C_0\cap C^\infty$
is bounded, and
\[
 \norm{v_0}1\leq C_aM,\quad
 \norm{v_0}2\leq C_a\norm\mu{H^{-1}},\quad
 \norm{\nu_0}1\leq M,\quad \norm{\sigma_0}\M\leq2M.
\]
All derivatives of $\nu_0$ belong to $L^1\cap L^\infty$, and every
positive-order derivative of $V_0$ belongs to $L^2\cap L^\infty$.
\end{lemma}
\begin{proof}
For a bounded ball $B$, $\sup_y\int_B|K(x-y)|\dd<\infty$; Tonelli gives the
almost-everywhere definition and local integrability. The kernel's divergence,
curl, and Fourier identities give the distributional assertions. These identities
extend from compactly supported to arbitrary finite measures by truncation in
total variation. Equivalently, $m$ is locally integrable at frequency zero and
integrable against Schwartz functions, so uniform convergence of the Fourier
transforms of the truncated measures passes the formula to the limit.

For radial $a$, circulation gives
\[
 (K*a)(x)=K(x)\int_{|y|<|x|}a(y)\,dy.
\]
Consequently $L_a=K-K*a$ vanishes outside $B_A$ and has only the integrable
$|x|^{-1}$ singularity at zero. Thus $L_a\in L^1$ and Young's inequality gives
the $L^1$ estimate for $v_0$. Evenness gives
$|1-\widehat a(\xi)|\leq C_a\min(|\xi|^2,1)$; hence
\[
 (1+|\xi|^2)|(1-\widehat a(\xi))m(\xi)|^2\leq C_a.
\]
Plancherel proves the $L^2$ estimate, using exactly the inhomogeneous hypothesis.
The multipliers $m\widehat a$ and $\xi^\alpha m\widehat a$ are integrable;
for $|\alpha|\geq1$ the latter multipliers $\xi^\alpha m\widehat a$ are also square integrable. Fourier inversion and the
Riemann--Lebesgue lemma give the assertions about $V_0$. The convolution
identities and remaining bounds follow by finite-measure convolution. In
particular $K*\mu=V_0+v_0\in L^2_{\rm loc}$.
\end{proof}

\paragraph{External input E1, used only here.}
Crippa--Stefani, \emph{An elementary proof of existence and uniqueness for the
Euler flow in localized Yudovich spaces}, Calc. Var. PDE 63, 168 (2024),
Theorem 3.3, with Assumptions 2.1 and 3.1 and Definition 3.2,
\url{https://link.springer.com/article/10.1007/s00526-024-02750-4}.
The specialization used here supplies a global Lagrangian vorticity solution for
arbitrary signed $L^1\cap L^\infty$ data on $\R^2$, with canonical Biot--Savart
velocity almost everywhere in time, bounded vorticity and velocity on finite
horizons, a log-Lipschitz velocity modulus, and nonincreasing $L^1$ norm.
Definition 3.2 includes the weak initial-value vorticity equation and the
Lagrangian pushforward representation almost everywhere in time. It does not require compact support or finite global energy.
The kernel hypotheses hold because $|K(z)|=(2\pi|z|)^{-1}$,
$|K(z)-K(w)|=|z-w|/(2\pi|z||w|)$, and $\Div K=0$; take $p=3$ in Assumption~3.1, whose bounded
$1/3$-H\"older velocity mapping is supplied by Theorem~2.2, and note that
bounded vorticity has uniformly bounded local $L^3$ norms.
The source is not being used as an unquoted smoothness theorem.

\begin{lemma}[Bounded background and its trace]\label{lem:background}
The datum $\nu_0$ in Lemma~\ref{lem:decomp} has a global bounded-vorticity Euler evolution
$\nu$, with velocity $V=K*\nu$ and a pressure $p_V$, such that on each finite
horizon $[0,T]$,
\[
 \sup_t(\norm{\nu(t)}1+\norm{\nu(t)}\infty+
 \norm{V(t)}\infty+\norm{V(t)}3)<\infty,
 \quad V(t)-V_0\in C([0,T];L^2).
\]
Moreover $\norm{V(t)-V_0}2\leq
t\norm V{L^\infty_{t,x}}\norm{\nu_0}2$.
\end{lemma}
\begin{proof}
Apply E1, using its $L^1$ and $L^\infty$ bounds and canonical velocity.
No smoothness upgrade is needed in this construction.
We first select an every-time representative.
The log-Lipschitz modulus is Osgood, so the flow and
its inverse are unique, continuous, and have bounded displacement on finite
horizons. Mollify the velocity only in space by a probability kernel.
The resulting smooth divergence-free flows preserve Lebesgue measure by
the Jacobian equation. Their uniform Osgood modulus and uniform convergence
of the velocities give convergence of flows and inverse flows on compact
time-space sets. For a compact continuous test, bounded displacement confines
its transported support to a fixed bounded set; dominated convergence passes
volume preservation to the original flow $X_t$.

Define $\nu(t)=\nu_0\circ X_t^{-1}$ at every time. It agrees with the source
solution almost everywhere in time and preserves the $L^1$, $L^2$, and
$L^\infty$ norms. Approximating $\nu_0$ in $L^1$ or $L^2$ by compact smooth
functions and using volume preservation and continuity of the inverse flow
gives $\nu\in C_tL^1\cap C_tL^2$. Set $V(t)=K*\nu(t)$ at every time.
It agrees with the source velocity almost everywhere, so its equation is
unchanged. For $f=\nu(t)-\nu(s)$ the near--far bound
$\norm{K*f}\infty\leq r\norm f\infty+(2\pi r)^{-1}\norm f1$
shows continuity of $V$ in $L^\infty$ by choosing $r$ small and then $t-s$
small. Thus this representative has the canonical initial velocity and
continuous positive-time local velocity pairings.

The Biot--Savart near--far split also bounds the velocity:
\[
 \norm V\infty\leq C\bigl(r\norm\nu\infty+r^{-1}\norm\nu1\bigr)
 \quad(r>0).
\]
Thus $V\nu$ belongs to both $L^1$ and $L^2$ uniformly on finite horizons.

For completeness, $K*\mu$ for any finite measure satisfies
\[
 |\{|K*\mu|>\lambda\}|\leq C\norm\mu\M^2\lambda^{-2}.
\]
Indeed split $K$ at radius $r=\norm\mu\M/(\pi\lambda)$; the far part has
supremum at most $\lambda/2$, and the near part has $L^1$ norm at most
$Cr\norm\mu\M$. Chebyshev proves the bound. Combining it with boundedness
of $V$ and layer cake gives $\norm V3^3\leq C\norm\nu1^2\norm V\infty$.

Integrate the vorticity equation:
\[
 \nu(t)-\nu_0=-\Div F_t,\qquad F_t=\int_0^t V(s)\nu(s)\,ds\in L^1\cap L^2.
\]
The Fourier multiplier of $-K*\Div$ has operator norm at most one, yielding
$V(t)-V_0=-K*\Div F_t$ and the stated estimate and continuity. Likewise
$\partial_tV,\nabla V\in L^\infty_tL^2$. The same argument between
any two times gives $\norm{V(t)-V(s)}2\leq C_T|t-s|$. With
$B_T=\sup_t\norm{V(t)}\infty$, interpolation gives
$\norm{V(t)-V(s)}3^3\leq2B_T\norm{V(t)-V(s)}2^2$.
Therefore $V\in C([0,T];L^3)$. The vector
$\partial_tV+(V\cdot\nabla)V$ is an $L^2$ curl-free field by transport;
it is a pressure gradient in distributions. Thus the velocity Euler equation
used below is indeed available without a global-energy assumption on $V$.
\end{proof}

\section{An initial Euler--Reynolds subsolution}

The convention for Euler--Reynolds is
\begin{equation}\label{eq:ER}
 \partial_tu+\Div(u\otimes u)+\nabla p=\Div R,
 \qquad \Div u=0,
\end{equation}
where $R$ is symmetric. Let $b=\mathcal F^{-1}(e^{-|\xi|^4})$ and
$b_t(x)=t^{-2}b(x/t)$ for $t>0$. It is a real radial Schwartz kernel of
integral one; it need not be positive. Write $C_b=\norm b1$.

\begin{proposition}[Initial Euler--Reynolds seed]\label{prop:seed}
For the preceding decomposition and background, put
\[
 W(t)=b_t*v_0,\qquad u^{(0)}=V+W,\qquad p^{(0)}=p_V,
\]
\[
 R^{(0)}=W\otimes W+V\otimes W+W\otimes V
 -4t^3\bigl[\nabla\Delta W+(\nabla\Delta W)^T\bigr].
\]
This is a positive-time distributional solution of \eqref{eq:ER}. On $0<t\leq T$,
\begin{align*}
 \norm{u^{(0)}-V}2&\leq\norm{v_0}2,\\
 \norm{\curl u^{(0)}}1&\leq\norm{\nu_0}1+C_b\norm{\sigma_0}\M,\\
 \norm{R^{(0)}}1&\leq\norm{v_0}2^2
 +2C_b\norm V\infty\norm{v_0}1+Ct\norm{\sigma_0}\M.
\end{align*}
It has strong $L^2_{\rm loc}$ initial trace $u_0$ and canonical velocity
$u^{(0)}(t)=K*(\nu(t)+b_t*\sigma_0)$.
On compact positive-time intervals $R^{(0)}\in C_tL^1$,
$u^{(0)}\in C_tL^3$, and the velocity is strongly continuous in local $L^2$.
\end{proposition}
\begin{proof}
The multiplier of $b_t$ is $e^{-t^4|\xi|^4}$, so
$\partial_tW=-4t^3\Delta^2W$. Since $\Div W=0$,
$\Div[\nabla\Delta W+(\nabla\Delta W)^T]=\Delta^2W$.
Expanding $(V+W)\otimes(V+W)$ proves \eqref{eq:ER} with the stated signs.
Plancherel gives the $L^2$ bound and $W(t)\to v_0$ strongly in $L^2$.
Finite-measure convolution gives $\curl W=b_t*\sigma_0$ and the curl bound.
Furthermore $\norm W1\leq C_b\norm{v_0}1$. The identity
$\Delta W=\nabla^\perp\curl W$ yields
\[
 \norm{\nabla\Delta W}1\leq C\norm{D^2b_t}1\norm{\sigma_0}\M
 =Ct^{-2}\norm{\sigma_0}\M.
\]
Thus the derivative stress costs $Ct\norm{\sigma_0}\M$, and the other
terms are bounded as displayed. Lemma~\ref{lem:background} proves the trace.
Fubini against compact tests, or the Fourier identity, gives the canonical
convolution. On positive-time compact intervals, convolution bounds give
$W\in C_t(L^1\cap L^2\cap L^\infty)$ and hence
$W\in C_tL^{3/2}\cap C_tL^3$. The derivative stress is continuous
in $L^1$ by continuity of $D^2b_t$ in $L^1$. Products $W\otimes W$
are $L^1$-continuous by $C_tL^2$, while $V\otimes W$ and its transpose
are $L^1$-continuous by $V\in C_tL^3$ and $W\in C_tL^{3/2}$.
This proves all asserted continuity without differentiating the old background.
\end{proof}

\begin{remark}[Why a subcritical global increment norm is required]\label{rem:cross}
Every repair below must have an \emph{arbitrarily small} $L^{3/2}$ velocity
increment, in addition to its $L^2$ and curl-variation estimates. Indeed,
$\norm{V\otimes w}1\leq\norm V3\norm w{3/2}$ closes the whole-plane
cross-stress bound. Finite energy and finite curl variation do not suffice:
a smooth radial swirl equal to $e_\theta/(r\log r)$ for $r\geq e^2$ belongs
to $L^2$, has integrable curl, but is not in $L^{3/2}$. Against a smooth
unit-circulation background equal to $e_\theta/(2\pi r)$ at infinity,
$\int|V\otimes w|=\infty$ because $\int_R^\infty dr/(r\log r)=\infty$.
Multiplying $w$ by a small constant does not repair this divergence.
\end{remark}

\section{Endpoint repair: elementary building blocks}

\subsection{An exact translating dipole}
The Bessel functions used below may be defined by their convergent power
series, whose termwise derivatives give the Bessel equation and
$J_1'(z)=J_0(z)-J_1(z)/z$.
There is a positive zero $\beta\in(3,4)$ of $J_1$.
For completeness, the alternating series gives
$J_1(3)\geq(3/2)(223/1024)>0$ by truncating after its cubic term in
$(3/2)^2$, and $J_1(4)\leq-2/45<0$ by truncating after its fourth such term.
The omitted terms alternate with decreasing magnitudes.
The intermediate value theorem gives the asserted zero.
If $J_0(\beta)=0$ also, the recurrence gives $J_1'(\beta)=0$.
Uniqueness for the regular second-order Bessel ODE near $\beta>0$ would then
make $J_1$ identically zero, contrary to $J_1(z)=z/2+O(z^3)$.
Thus $J_0(\beta)\ne0$.
In polar coordinates set
\[
 \psi(\rho,\theta)=
 \begin{cases}
 -\rho\sin\theta+\dfrac{2J_1(\beta\rho)}{\beta J_0(\beta)}\sin\theta,
 &\rho\leq1,\\
 -\sin\theta/\rho,&\rho\geq1,
 \end{cases}
 \qquad D=\nabla^\perp\psi.
\]
With $q=\psi+x_2$, $D-e_1=\nabla^\perp q$ and
$\curl D=-\beta^2q$ in $B_1$, while $\curl D=0$ outside.
The relations $J_1(\beta)=0$ and $J_1'(\beta)=J_0(\beta)$ give the matching
values and first derivatives of the stream function.
The Bessel equation also gives $J_1''(\beta)=-J_0(\beta)/\beta$,
so the second radial derivatives match.
Tangential differentiation of the matching traces gives the other second
derivatives; hence the velocity is continuously differentiable across the circle.
In particular $q=0$ on
the circle and $(D-e_1)\cdot\nabla\curl D=0$ distributionally. The pressure
\[
 P=-\tfrac12|D-e_1|^2-\tfrac{\beta^2}{2}q^2\mathbf1_{B_1}
\]
gives $-\partial_1D+\Div(D\otimes D)+\nabla P=0$. This follows from
$(w\cdot\nabla)w=(\curl w)w^\perp+\nabla(|w|^2/2)$, and matching on the
circle leaves no boundary distribution. Outside the unit disk,
$D=\nabla(-x_1/|x|^2)$; $D$ is even under $x\mapsto-x$.

For a unit direction $\xi\in\mathbb S^1$, rotate the velocity and pressure profiles together, and write $D_r(x)=r^{-1}D(x/r)$ and $P_r(x)=r^{-2}P(x/r)$.
The direction is suppressed where no confusion arises. Then
\begin{equation}\label{eq:dipole}
 -r^{-1}\partial_\xi D_r+\Div(D_r\otimes D_r)+\nabla P_r=0,
 \quad \Div D_r=0,
\end{equation}
and
\begin{equation}\label{eq:dipole-norms}
 \norm{D_r}2=C_2,\quad \norm{\curl D_r}1=C_\omega,\quad
 \norm{D_r}{3/2}=C_{3/2}r^{1/3},\quad \norm{\nabla D_r}1=C_\nabla.
\end{equation}
The constants are finite by boundedness in the core and the exterior
$|x|^{-2}$ decay; the curl is supported in $B_r$.

\subsection{A symmetric globally integrable anti-divergence}
For a vector field $f$ supported in $B_R$, with
\[
 \int f=0,\qquad \int x_jf_i(x)\dd=0\quad\hbox{for every }i,j,
\]
define
\[
 (\mathcal Af)_{ij}=\partial_i\Delta^{-1}f_j+
 \partial_j\Delta^{-1}f_i-\delta_{ij}\Div\Delta^{-1}f.
\]
It is symmetric and $\Div\mathcal Af=f$. Its kernel has degree $-1$.
On $B_{3R}$, integration of that kernel gives $CR\norm f1$. Outside $B_{3R}$,
subtract its constant and linear Taylor terms using the two moment conditions;
the remainder is at most $CR^2|x|^{-3}\norm f1$. Integration gives
\begin{equation}\label{eq:anti}
 \norm{\mathcal Af}1\leq CR\norm f1.
\end{equation}
The output need not be compactly supported. Cancellation of only the mean
would leave a generally nonintegrable $|x|^{-2}$ far field.

Differentiation of the profile gives
\[
 \partial_rD_r=r^{-1}D_r-\Div(D_r\otimes x/r).
\]
The vector $f_r=\Div(D_r\otimes x/r)$ is supported in $B_r$, because the
exterior dipole is homogeneous of degree $-2$; $\norm{f_r}1\leq C$.
Evenness gives all first moments zero. Its mean is zero by the boundary flux
on any circle outside $B_r$, since the angular mean of the exterior dipole
is zero. Therefore $F_r=-\mathcal Af_r$ satisfies
\begin{equation}\label{eq:radial}
 \partial_rD_r=r^{-1}D_r+\Div F_r,
 \qquad \norm{F_r}1\leq Cr.
\end{equation}

\subsection{Variable radius and a compact source}
On each visit take $r,\eta,X$ to be $C^1$ in time, with $r(t)>0$, $\eta(t)\geq0$, and $X'=\eta\xi/r$, and write $T_X f(x)=f(x-X)$. By
\eqref{eq:dipole}--\eqref{eq:radial}, the packet
$v(t,x)=\eta(t)D_{r(t)}(x-X(t))$ satisfies, after adding gradients to pressure,
\begin{equation}\label{eq:packet}
 \partial_tv+\Div(v\otimes v)+\nabla p
 =(\eta r)'T_XD_r/r+\Div(\eta r'T_XF_r).
\end{equation}
Choose a radial cutoff $\chi$ equal to zero on $B_1$ and one outside $B_2$,
let $\alpha=1/5$, and put
\[
 \Phi_\xi(x)=-\xi\cdot x/|x|^2,\quad
 \Pi_r(x)=r\chi(x/r^\alpha)\Phi_\xi(x),\quad
 Z_r=D_r-\nabla\Pi_r.
\]
For small $r$, $Z_r$ is even and supported in $B_{2r^\alpha}$; direct radial
integration gives $\norm{Z_r}1\leq Cr(1+|\log r|)$. Its integral is
$2\pi r\xi$: on a large circle the respective boundary integrals for $D_r$
and $\nabla\Pi_r$ are $\pi r\xi$ and $-\pi r\xi$.
Thus $S_r=Z_r/r$ has integral $2\pi\xi$ and all first moments zero.

Fix a nonnegative radial smooth probability bump $U\in C_c^\infty(B_1)$ and define $U_H(x)=H^{-2}U(x/H)$, supported in $B_H$.
If $2r^\alpha\leq H$, the difference $S_r-2\pi U_H\xi$ has its zeroth and
first moments zero. Equation~\eqref{eq:anti} gives
\[
 \norm{\mathcal A(S_r-2\pi U_H\xi)}1\leq CH(1+|\log r|).
\]
Equation~\eqref{eq:packet} therefore becomes
\begin{equation}\label{eq:packet-source}
 \partial_tv+\Div(v\otimes v)+\nabla p=\mathcal U+\Div E,
 \qquad \mathcal U=2\pi(\eta r)'U_H(x-X)\xi,
\end{equation}
where
\begin{equation}\label{eq:E-bound}
 E=\eta r'T_XF_r+(\eta r)'T_X\mathcal A(S_r-2\pi U_H\xi),\quad
 \norm E1\leq C|\eta r'|r+C|(\eta r)'|H(1+|\log r|).
\end{equation}
The endpoint estimate $\norm{\curl v}1=|\eta|C_\omega$ is independent of $r$.

\section{Two-scale raster and local stress repair}\label{sec:repair-interface}

\subsection{Stress coefficients and a fixed cycle}
For a prescribed symmetric stress piece $S$, use the directions
$e_1,e_2,(e_1+e_2)/\sqrt2,(e_1-e_2)/\sqrt2$.
For a fixed smooth compactly supported tensor $S$, choose a nonnegative
compact smooth function $h$ equal to one on a neighborhood of its support.
For $\gamma>0$ put
\[
 A=\sqrt{S_{12}^2+\gamma^2h^2},\qquad
 B=\sqrt{S_{11}^2+S_{22}^2+\gamma^2h^2}.
\]
These functions are smooth: near the zero set of $h$, the entries of $S$
already vanish, so both square roots reduce to $\gamma h$.
They satisfy $A\geq|S_{12}|$ and $B\geq\max(S_{11},S_{22})$.
Define
\[
 a_1=B-S_{11},\quad a_2=B-S_{22},\quad
 a_3=A-S_{12},\quad a_4=A+S_{12},\quad q=A+B
\]
Then $-S+qI=\sum_i a_i\xi_i\otimes\xi_i$.
All coefficients are smooth, compactly supported, and nonnegative, with
$\sum_i\norm{a_i}1\leq C\norm S1+C\gamma\norm h1$.
The final term can be made arbitrarily small before choosing fine scales.
The discarded spatial tail remains a small residual. On a compact time
interval, continuity into $L^1$ makes this tail truncation uniform: the
continuous image is compact in $L^1$, and a finite covering reduces the claim
to finitely many fixed integrable functions.

Fix a coefficient and its rectangle in coordinates $(s,y)$ adapted to $\xi$.
Write $a$ for the original compactly supported coefficient.
On the scanning rectangle define $c=a+e$ with a strictly positive floor
$e>0$, and use constant-$e$ longitudinal buffers where $a=0$.
Write $\varepsilon_{\rm floor}$ for the total added floor and buffer mass,
including the fixed factors from the ramp durations below.
This mass is chosen arbitrarily small.
All rectangle, floor, and derivative constants are fixed before the fine
parameter tends to zero. We use two distinct scales
\begin{equation}\label{eq:scales}
 d=H^2\quad\hbox{(row spacing)},\qquad
 \tau\asymp H^6\quad\hbox{(cycle duration)}.
\end{equation}
The auxiliary bump has width $H$, not $d$. One packet visits rows $y=kd$
sequentially, with no simultaneous packets. On the interior of a row,
$\eta$ is constant and $r(s)=r_0c(s,kd)$. Thus
\[
 dt=(r_0/\eta)c(s,kd)\,ds,
 \qquad\int\eta^2U_H(x-X(t))\,dt
 =\eta r_0\int c(s,kd)U_H(x-s\xi-kd\xi^\perp)\,ds.
\]
At entry use a fixed smooth temporal ramp $\zeta:[0,1]\to[0,1]$,
equal to zero to all orders at zero and equal to one with all positive-order
derivatives zero at one.
At exit use the reversed ramp $\zeta(1-\cdot)$.
Write $I_\zeta=\int_0^1\zeta>0$.
Throughout either buffer $r=r_0e$ is constant.
To traverse a
buffer of length $L$, choose ramp duration
$r_0eL/(\eta I_\zeta)$.
Let $B_i$ be the sum over rows of the interior line
integrals of $c_i=a_i+e_i$, plus $2e_iL/I_\zeta$ per row for the two ramps.
Allocate $\tau_i=\tau/4$ to direction $i$, and set
\begin{equation}\label{eq:normalization}
 \eta_i^2=\frac{\tau dB_i}{2\pi\tau_i},\qquad
 r_{0,i}=\frac{\tau d}{2\pi\eta_i}.
\end{equation}
The full scan takes exactly $(r_{0,i}/\eta_i)B_i=\tau_i$, and its coefficient
weight is exactly $2\pi\eta_ir_{0,i}/\tau=d$. With sufficiently fine $d$,
$dB_i\leq C\norm{a_i}1+C\varepsilon_{\rm floor}+o_H(1)$.
Writing $\varepsilon_{\rm reg}$ for the earlier regularization mass gives
$\eta_i\leq C\sqrt\delta+
C\sqrt{\varepsilon_{\rm floor}+\varepsilon_{\rm reg}}+o_H(1)$
when $\norm S1\leq\delta$.
If a direction
has zero mass it can be omitted; a positive floor otherwise makes
\eqref{eq:normalization} unambiguous.

On a ramp $|(\eta r)'|\leq C\eta_i^2$; on an interior row
$|\eta r'|=|(\eta r)'|\leq\eta_i^2\norm{\partial_\xi\log c_i}\infty$.
Constants are fixed before $H\downarrow0$. Accordingly $r_0\asymp H^8$,
all radii are bounded above and below by fixed multiples of $H^8$, and
$r^{1/5}=O(H^{8/5})=o(H)$. The intrinsic error
\eqref{eq:E-bound} is $O(H|\log H|)+O(H^8)$.

Integration by parts on a visit, using $\eta=0$ at its endpoints, gives
\begin{equation}\label{eq:average}
 \overline{\mathcal U}=\frac1\tau\int_{\rm cycle}\mathcal U\,dt
 =\Div\left(\frac{2\pi}{\tau}\int_{\rm cycle}
 \eta(t)^2U_H(x-X(t))\xi\otimes\xi\,dt\right).
\end{equation}
Indeed $\frac d{dt}(\eta rU_H(x-X))=(\eta r)'U_H-\eta^2\xi\cdot\nabla U_H$.
Write the tensor in parentheses, summed over four directions, as $A_{\rm av}$.
The raster quadrature obeys
\begin{equation}\label{eq:quad}
 \left\|d\sum_k\int a(s,kd)U_H(\,\cdot-s\xi-kd\xi^\perp)\,ds-a\right\|_1
 \leq C\left[d\bigl(\norm{\partial_\perp a}1+H^{-1}\norm a1\bigr)
 +H\norm{\nabla a}1\right].
\end{equation}
Here is the explicit Banach-valued rectangle-rule proof.
Put $F(y)=\int a(s,y)U_H(\cdot-s\xi-y\xi^\perp)\,ds$, as an $L^1_x$-valued
function.
Differentiating it gives
$\int\norm{F'(y)}1dy\leq\norm{\partial_\perp a}1+
\norm a1\norm{\partial_\perp U_H}1$.
On each row interval of length $d$, the fundamental theorem of calculus
bounds the integrated discrepancy from its endpoint sample by
$d$ times the integral of $\norm{F'}1$ on that interval.
Summing yields the first two terms in \eqref{eq:quad}.
The translation estimate
$\norm{a*U_H-a}1\leq CH\norm{\nabla a}1$ gives the third.
Estimate \eqref{eq:quad} applies to the original compact coefficient $a$.
The added floor and the ramp-weighted buffers contribute at most
$C\varepsilon_{\rm floor}+o_H(1)$ separately.
Since $d=H^2$, the quadrature error tends to zero.

\subsection{Temporal correction and the stress identity}
On each cycle $[t_k,t_k+\tau]$ put
\[
 Q(t)=-\mathbb P\int_{t_k}^t
 (\mathcal U(s)-\overline{\mathcal U})\,ds,
\]
The average here is taken over the entire four-direction cycle, not
separately on each direction slot. This ensures that its tensor is the
sum of all four directional averages throughout the cycle.
Here $\mathbb P=I-\nabla\Delta^{-1}\Div$ is used only on translates of
the fixed bump $U_H\xi$ and their time integrals. Its needed bounds have
a direct fixed-profile proof. Let $G=(2\pi)^{-1}\log|x|$ and
\[
 Z_{H,\xi}=U_H\xi-(D^2G*U_H)\xi.
\]
It has zero divergence and $\curl Z_{H,\xi}=\curl(U_H\xi)$.
At $H=1$ it is smooth and bounded in the core: write the Hessian convolution
as $\partial_iG*\partial_jU_1$, whose $|x|^{-1}$ singularity is integrable.
For $|x|>1$, radial averaging gives $G*U_1=G$, so
$Z_{1,\xi}=-D^2G\xi=O(|x|^{-2})$. Hence $Z_{1,\xi}\in L^p$ for every
$p>1$. Scaling gives
\[
 \norm{Z_{H,\xi}}p\leq C_p H^{2/p-2}.
\]
Minkowski applies to its translates in the time primitive. No general
Calderón--Zygmund $L^p$ theorem is needed for this use of $\mathbb P$.
For an explicit pressure sign, write
$q_{H,\xi}=\xi\cdot\nabla(G*U_H)$ and let
$\pi_{\mathcal U}(t)=2\pi(\eta r)'q_{H,\xi}(\cdot-X(t))$
for the active direction, with zero value when no packet is active.
Then $\mathcal U-\mathbb P\mathcal U=\nabla\pi_{\mathcal U}$.
Taking
$p_Q=-\pi_{\mathcal U}+\tau^{-1}\int_{\rm cycle}\pi_{\mathcal U}(s)\,ds$
gives the exact identity
$\partial_tQ+\mathcal U+\nabla p_Q=\overline{\mathcal U}$.
If $B$ bounds $|(\eta r)'|$, bump scaling gives
\begin{equation}\label{eq:Q}
 \norm Q2\leq C\tau H^{-1}B,
 \quad\norm Q{3/2}\leq C\tau H^{-2/3}B,
 \quad\norm{\curl Q}1\leq C\tau H^{-1}B.
\end{equation}
Here both $\mathcal U$ and its cycle average obey the needed norm bound.
The curl estimate follows directly from $\norm{\nabla U_H}1=O(H^{-1})$.
No $L^1$ boundedness of $\mathbb P$ is invoked. The corrector vanishes at both
cycle endpoints. Its first time derivative may jump there, so the new velocity correction
is continuous and, after the fixed spatial mollification below, piecewise
smooth in time; it is not asserted to
be globally time-smooth.

Given an old distributional Euler--Reynolds seed $(u,p,R)$, set $\widetilde u=u+v+Q$.
Combining \eqref{eq:packet-source}, \eqref{eq:average}, and the equation for
$Q$ gives, modulo a gradient,
\[
 \partial_t\widetilde u+\Div(\widetilde u\otimes\widetilde u)
 =\Div\widetilde R,
\]
\begin{align*}
 \widetilde R={}&R+A_{\rm av}-qI+E
 +u\otimes v+v\otimes u+u\otimes Q+Q\otimes u\\
 &+v\otimes Q+Q\otimes v+Q\otimes Q.
\end{align*}
The scalar $q$ is absorbed into pressure. Thus $\widetilde R=R-S+E_{\rm total}$,
where the averaging error is $S+A_{\rm av}-qI$ and all remaining terms are
shown explicitly. The displayed expression for $\widetilde R$ is symmetric, since each mixed product occurs with its transpose;
the cross terms satisfy
\begin{align*}
 \norm{u\otimes v}1&\leq\norm u3\norm v{3/2}
 \leq C\norm u3\eta r^{1/3},&
 \norm{u\otimes Q}1&\leq\norm u3\norm Q{3/2},\\
 \norm{v\otimes Q}1&\leq\norm v2\norm Q2,&
 \norm{Q\otimes Q}1&=\norm Q2^2.
\end{align*}
These terms tend to zero at fixed coarse seed as $H\downarrow0$.
The same estimates give an arbitrarily small $L^{3/2}$ increment.
A corrector on one direction slot alone would have the wrong instantaneous
average stress and is not used.

\subsection{Spatial smoothing, without remollifying the old seed}
The raw profile can be made spatially smooth on each fixed cycle by convolving
the \emph{new packet equation} with a fixed even compact smooth mollifier
$\rho_\ell$. Keep $\ell$ constant throughout that cycle and take it below its
positive minimum radius. Replace $v,Q,\mathcal U,E,A_{\rm av}$ by their
convolutions consistently. The packet equation then gains the symmetric
commutator
\[
 C_\ell=v_\ell\otimes v_\ell-(v\otimes v)_\ell,
\]
with the sign determined by moving it to the right-hand side as above.
Its $L^1$ norm tends uniformly to zero on the fixed cycle: the map
$t\mapsto v(t)$ is continuous into $L^2$, its image is compact there, and
approximation of the identity is uniform on compact subsets of $L^2$.
Norm bounds do not grow under probability convolution; evenness and the two
moment cancellations are preserved. The constant smoothing scale eliminates
any additional time derivative of the convolution kernel. Zero endpoint
values of $v$ and $Q$ are preserved. The old seed is never remollified, so no
new velocity jump is created at a pre-existing boundary.

\begin{lemma}[Prescribed-piece local endpoint repair, LER]\label{ass:LER}
Consider the following admissible class: positive-time Euler--Reynolds
velocities are continuous into local $L^2$ and global $L^3$.
Stresses are symmetric and locally bounded in $L^\infty_tL^1_x$.
Their restrictions to the open pieces of a locally finite positive-time
partition admit continuous $L^1$ extensions to each piece's closure.
The two one-sided extensions need not agree at a boundary.
Actual stress values at those discrete times are chosen separately subject
to the same norm bound; they do not enter the distributional equation.
The Euler--Reynolds equation holds in distributions. No derivative bound
on the old velocity is required by the local construction.

Let $(u,R)$ be any current pair in this class, $J\Subset(0,\infty)$ an
open interval, and $S$ a prescribed symmetric tensor compactly supported in
time in $J$, piecewise continuous into $L^1$ on the inherited partition.
Write $\delta=\sup_J\norm S1$.
For every $\varepsilon,\zeta,\kappa,r_*>0$
there are $w,E_{\rm total}$ compactly supported in time in $J$ such that
\[
 (\widetilde u,\widetilde R)=(u+w,R-S+E_{\rm total})
\]
is again an Euler--Reynolds pair in the admissible class,
$\Div w=0$, $E_{\rm total}$ is symmetric, and
\[
 \sup_J\norm{E_{\rm total}}1\leq\varepsilon,\qquad
 \sup_J\bigl(\norm w2+\norm{\curl w}1\bigr)\leq C\sqrt\delta+\zeta,
 \qquad \sup_J\norm w{3/2}\leq\kappa.
\]
The constant $C$ is independent of the seed derivatives, support, patch
location, and patch length; fine parameters may depend on these quantities
and on the finite old $L^3$ norm. The new correction is continuous into
$L^2\cap L^3\cap L^{3/2}$, and only finitely many new time boundaries are
introduced in a compact patch. Old boundaries are included in the new
partition.
Every principal radius can be required to be at most $r_*$.
All outputs paste within the admissible class.
\end{lemma}
\begin{proof}
If $\delta=0$, take $w=E_{\rm total}=0$.
Suppose henceforth that $\delta>0$.
Choose $[a,b]\Subset J$ containing the temporal support of $S$ in its interior.
Only finitely many inherited partition boundaries meet this interval.
On each closed one-sided continuity piece, uniform continuity into $L^1$
permits a finite coarse subdivision such that on each coarse cell $I_k$
there is a fixed symmetric tensor $S_k\in C_c^\infty$ satisfying
\begin{equation}\label{eq:coarse-approx}
 \sup_{t\in I_k}\norm{S(t)-S_k}1\leq\varepsilon_0,\qquad
 \norm{S_k}1\leq\delta+\varepsilon_0.
\end{equation}
In this display, $S$ on a closed $I_k$ means the one-sided continuous
extension from its interior.
No single actual boundary value is required to equal two different traces.
First freeze $S$ at one time with error at most $\varepsilon_0/2$, then approximate
that single $L^1$ tensor by a compact smooth symmetric tensor with error
at most $\varepsilon_0/2$.
There are only finitely many coarse cells.
Thus their chosen supports, spatial derivatives, and decomposition parameters
have finite maxima.
No derivative of the actual prescribed $S$ is taken.

On each $I_k$, decompose the constant tensor $-S_k$ modulo $q_kI$
as above.
Choose its regularization parameter and the positive floors so that all
added coefficient mass is arbitrarily small.
The rectangular buffers lie outside the support of the original coefficient,
where its value is zero before adding the floor.
The constant-floor buffer errors are bounded by a fixed multiple of their
total added mass.
In particular, shrinking those masses controls the averaging error and the
additive amplitude slack independently of the size of the rectangle.

Fix the coarse tensor, rectangles, floors, and ramp functions.
Divide $I_k$ into $N_k$ equal complete four-direction cycles.
Put $\tau_k=|I_k|/N_k$, $H_k=\tau_k^{1/6}$, and $d_k=H_k^2$.
Consequently $\tau_k=H_k^6$ holds exactly, rather than only asymptotically,
and every endpoint of $I_k$ is a complete-cycle endpoint.
Use the exact normalization \eqref{eq:normalization} on each cycle.
At fixed positive floors, the quantities $d_kB_i$ converge to positive finite
limits as $N_k\to\infty$.
The amplitudes therefore have positive lower and finite upper bounds
independent of large $N_k$, with bounds permitted to depend on the fixed floors.
The leading upper bound is
\[
 \eta_i\leq C\sqrt{\delta+\varepsilon_0}
   +C\sqrt{\hbox{added mass and vanishing row-rule errors}}.
\]
The constant multiplying $\sqrt{\delta+\varepsilon_0}$ depends only on the fixed
four directions and dipole, not on derivative norms or the rectangle.

All radii on the fixed coarse cell are bounded above and below by positive
constant multiples of $H_k^8$.
The constants may depend on its coefficients and floors.
Thus $2r^{1/5}\leq H_k$ and $r\leq r_*$ hold for all sufficiently large $N_k$.
The quantities $|(\eta r)'|$ and $|\eta r'|$ stay bounded independently
of $N_k$, because their formulas involve only the uniformly bounded peak amplitudes, fixed
ramp functions, and the fixed logarithmic derivatives of the positive coefficients.
Equation \eqref{eq:E-bound} makes the intrinsic stress
$O(H_k|\log H_k|)+O(H_k^8)$.
Equation \eqref{eq:Q} makes the corrector $O(H_k^5)$ in $L^2$ and
curl $L^1$, and $O(H_k^{16/3})$ in $L^{3/2}$.
The principal $L^{3/2}$ norm is $O(H_k^{8/3})$ at fixed coarse parameters.
Equation \eqref{eq:quad} makes the quadrature error tend to zero.
Every cross term with the actual old $u$ tends to zero by the displayed
Hölder estimates and the finite bound $\sup_{[a,b]}\norm u3$.
There is no time-freezing derivative estimate here: the constructed tensor
$S_k$ is constant on the coarse cell, and the already charged discrepancy
is precisely \eqref{eq:coarse-approx}.

Finally choose a fixed spatial mollification scale on each complete cycle.
The finite set of cycles has a positive minimum radius.
The commutator can therefore be made smaller than any remaining stress budget,
while convolution decreases all required velocity and curl norms.
It changes the averaged tensor to its convolution.
The resulting additional averaging error also tends to zero in $L^1$ as
the mollification scale tends to zero, since that tensor is integrable.
No old field is mollified.

The exact stress identity on $I_k$ is
$\widetilde R=R-S_k+E_k^{\rm loc}$.
It follows by adding the old equation to the new packet and corrector
identities, taking the average over the whole four-direction cycle, and
subtracting $q_k$ from the intermediate pressure.
At times strictly inside a coarse cell, set
\[
 E_{\rm total}(t)=S(t)-S_k+E_k^{\rm loc}(t)\qquad(t\in I_k).
\]
Then $\widetilde R=R-S+E_{\rm total}$ exactly.
Choose $\varepsilon_0\leq\varepsilon/2$ and all local stress errors at most
$\varepsilon/2$.
This proves the requested stress bound.

Choose $\varepsilon_0$ also so that $C\sqrt\varepsilon_0\leq\zeta/3$.
Next choose added coefficient masses small enough for their amplitude cost
to be at most $\zeta/3$.
Finally choose the cycle counts large enough that the row-rule amplitude
error and the $L^2$ and curl cost of $Q$ use at most the remaining $\zeta/3$.
Using $\sqrt{\delta+\varepsilon_0}\leq\sqrt\delta+\sqrt\varepsilon_0$ proves
the independent leading increment estimate.
The same cycle counts can make the two $L^{3/2}$ contributions sum to at
most $\kappa$ and all radii at most $r_*$.
All these requirements have already been shown to improve as the counts
increase at fixed coarse data, so this is one compatible parameter choice.

Every principal packet and $Q$ vanishes at each completed cycle endpoint.
Extend $w$ by zero outside $[a,b]$.
Its continuous zero traces exclude temporal delta functions.
Define $E_{\rm total}=0$ outside this interval; the prescribed $S$ is zero there.
On the finite set of inherited, coarse-cell, and full-cycle boundary times,
choose $E_{\rm total}(t)=0$ and
$\widetilde R(t)=R(t)-S(t)$.
These choices maintain the exact algebra and the pointwise error bound,
and do not change the distributional equation or one-sided extensions.
The error stress may jump at $a,b$ and at cycle boundaries, but spatial
divergence creates no temporal delta from those jumps.
Both output supports are compact subsets of $J$.

The principal packet is continuous in global $L^2$ and $L^{3/2}$ by
positive-radius scaling and translation continuity, including its zero limits
at resets.
Fixed spatial smoothing makes it continuous into $L^3$ by Young's inequality
$\norm{\rho_\ell*v}3\leq C\ell^{-1/3}\norm v2$.
The projected-bump time integral makes $Q$ continuous into all three spaces.
For each fixed construction these bounds are finite, without asserting a
uniform $L^3$ bound through the infinite iteration.

Every local intrinsic stress is continuous into $L^1$ on its closed cycle.
For the anti-divergence terms this follows from \eqref{eq:anti} on a common
support ball and the $L^1$ continuity of the compact zero-moment inputs.
Translations and positive-radius rescalings preserve the asserted continuity.
At packet resets the prefactors in those intrinsic stresses tend to zero.
The commutator and quadratic terms are continuous by their $L^2$ products.
The old--new terms are continuous because
\[
 \norm{u(t)\otimes w(t)-u(s)\otimes w(s)}1
 \leq\norm{u(t)-u(s)}3\norm{w(t)}{3/2}
 +\norm{u(s)}3\norm{w(t)-w(s)}{3/2}.
\]
Adding $S-S_k$ leaves the error piecewise continuous into $L^1$.
Thus the new stress and velocity remain in the stated admissible class.
Only finitely many coarse cells and cycles were used in $[a,b]$.
\end{proof}

\section{Global time iteration and exact-target closure}

\begin{lemma}[Zero initial trace from anchor times]\label{lem:anchors}
Suppose $(U,R)$ and $(\widetilde U,\widetilde R)$ are positive-time
Euler--Reynolds pairs with continuous local $L^2$ velocities, uniformly
bounded local kinetic energies and global $L^1$ stresses on $(0,T)$.
Suppose $W=\widetilde U-U$ is uniformly bounded in global $L^2$ and
$W(t_k)=0$ for a sequence $t_k\downarrow0$. Then $W(t)\rightharpoonup0$
in global $L^2$ as $t\downarrow0$.
\end{lemma}
\begin{proof}
For a compact smooth divergence-free spatial test $\psi$, subtract the
Euler--Reynolds equations:
\[
 \frac d{dt}\int W\cdot\psi
 =\int(\widetilde U\otimes\widetilde U-U\otimes U):\nabla\psi
 -\int(\widetilde R-R):\nabla\psi.
\]
The right side is uniformly bounded. The pairing is Lipschitz, including
across piecewise-time boundaries by the distributional equation and
continuity, and its zero values at $t_k$ force its limit to zero.
For any sequence $t_n\downarrow0$, the bounded sequence $W(t_n)$ has a weakly convergent $L^2$ subsequence. Its limit
is divergence free and annihilates all such $\psi$; testing
$\psi=\nabla^\perp\phi$ also makes its curl zero. Fourier variables force
this global $L^2$ vector field to vanish: away from frequency zero it is
both parallel and perpendicular to $\xi$, and an $L^2$ function has no
zero-frequency point mass. All weak cluster points are zero, proving the claim.
\end{proof}

\begin{proposition}[Global iteration]\label{prop:iteration}
The seed in Proposition~\ref{prop:seed} can be repaired
on $(0,\infty)$ to a global weak Euler solution with exact initial trace,
locally uniformly finite curl variation, and a globally $L^2$ perturbation
of the seed at each time, bounded uniformly on finite time horizons.
\end{proposition}
\begin{proof}
Choose a nonnegative $\eta\in C_c^\infty((1/2,2))$ equal to one on $[3/4,3/2]$, whose positive dyadic dilates
cover $(0,\infty)$ and set
\[
 I_j=(2^{j-1},2^{j+1}),\qquad
 \chi_j(t)=\frac{\eta(t/2^j)}{\sum_{k\in\mathbb Z}\eta(t/2^k)}.
\]
This is a locally finite partition of unity with
$\supp\chi_j\Subset I_j$. Color patches by $j\bmod3$.
Within each color the intervals are disjoint and leave gaps; for a given
$j$ of that color, $3\cdot2^j$ lies between $I_j$ and $I_{j+3}$.
These gap times tend to zero as $j\to-\infty$. Use $N=3$ as an overlap bound.

Choose a smooth nondecreasing majorant $\Theta\geq1$, finite on every finite
horizon including zero, such that $\norm{R^{(0)}(t)}1\leq\Theta(t)$.
It exists from Proposition~\ref{prop:seed}. Let $\delta_q=16^{-q}$ and
inductively require $\norm{R^{(q)}(t)}1\leq\delta_q\Theta(t)$.
Freeze the prescribed pieces $S_{qj}=\chi_jR^{(q)}$ before any repair at
stage $q$. Set $M_j=\sup_{I_j}\Theta$, $m_j=\inf_{I_j}\Theta\geq1$.
For patch $j$ apply LER with
\[
 \varepsilon_{qj}=\delta_{q+1}m_j/N,\qquad
 \zeta_{qj}=\sqrt{\delta_q}/N,\qquad
 \kappa_{qj}=2^{-q}/N,\qquad r_*=1.
\]
Process the three colors in order. Within one color disjoint supports
allow simultaneous pasting, with no cross term between different patches.
Between colors use the current repaired pair as background, while still
removing the originally frozen piece $S_{qj}$. This prescribed-piece
feature is why LER is stated with $S$ independently of the current $R$.
Each compact positive-time interval meets only finitely many patches,
and LER introduces only finitely many new boundaries per patch.
Thus the finite-stage admissible class is retained.

At the end of the stage,
\[
 u^{(q+1)}=u^{(q)}+w_q,\quad w_q=\sum_jw_{qj},\qquad
 R^{(q+1)}=\sum_jE_{qj}.
\]
There are at most $N$ active summands at any time, and each has
$m_j\leq\Theta(t)$. Therefore
$\norm{R^{(q+1)}(t)}1\leq\delta_{q+1}\Theta(t)$.
During a partially completed color stage the stress is
\[
 \left(1-\sum_{j\ \mathrm{processed}}\chi_j\right)R^{(q)}
 +\sum_{j\ \mathrm{processed}}E_{qj},
\]
whose $L^1$ norm is at most $(\delta_q+\delta_{q+1})\Theta(t)$.
If $0<t\leq T$ and a patch is active, $I_j\subset(0,4T)$, so
\begin{equation}\label{eq:increment-stage}
 \sup_{0<t\leq T}\left(\norm{w_q(t)}2+
 \norm{\curl w_q(t)}1\right)
 \leq[NC\sqrt{\Theta(4T)}+1]\sqrt{\delta_q}.
\end{equation}
These bounds are summable in $q$ and also bound intermediate color increments.
The arbitrary $L^{3/2}$ budgets give summable subcritical increments; their
more immediate role is making each old--new global cross stress as small
as needed in LER. No uniform-in-stage $L^3$ norm is asserted.

For each color the total correction vanishes at its gap times approaching
zero. Both old and new stresses and local energies are uniformly bounded
near zero, and the global $L^2$ correction is bounded by
\eqref{eq:increment-stage}. Lemma~\ref{lem:anchors} shows it has weakly zero
initial trace. There are only three colors. Induction therefore preserves
exact initial trace at every stage without requiring small initial energy
or a nearby-endpoint replacement.

The correction series converges uniformly in global $L^2$ on $(0,T]$
for each finite $T$. At positive times the finite stages are locally
strongly $L^2$ continuous, so their locally uniform limit is too. At zero,
split the series into finitely many terms and a uniformly small $L^2$ tail;
each finite correction has weakly zero trace by the previous paragraph.
Together with the seed trace this gives
$u\in C_w([0,\infty);L^2_{\rm loc})$ and $u(0)=u_0$.
At each fixed $t>0$, the curl series converges absolutely in $L^1$ by
\eqref{eq:increment-stage}, and its sum is the distributional curl of the
velocity limit. Consequently
\[
 \sup_{0<t\leq T}\norm{\curl u(t)}1\leq
 \sup_{0<t\leq T}\norm{\curl u^{(0)}(t)}1+
 [NC\sqrt{\Theta(4T)}+1]\sum_q\sqrt{\delta_q}<\infty.
\]
Add the initial measure norm to obtain the bound including zero.
Divergence zero is preserved distributionally.

Strong $L^2$ convergence on compact positive-time cylinders passes the
quadratic term to the limit in $L^1$, while
$\norm{R^{(q)}(t)}1\leq\delta_q\Theta(t)$ removes the stress. Thus Euler
holds on $(0,\infty)$. For a test meeting zero, use a temporal cutoff rising
on $(\varepsilon,2\varepsilon)$. Uniform local energy makes the omitted
convection integral tend to zero. Weak local $L^2$ continuity identifies
the cutoff-derivative term with $\int u_0\cdot\varphi(0)$; replacing
$\varphi(t)$ by $\varphi(0)$ in this term incurs an $O(\varepsilon)$ error.
This proves precisely the target's initial-value weak formulation.
\end{proof}

\begin{lemma}[Canonical normalization from a fixed background]\label{lem:canonical}
Suppose $u=U+v$, where $U\in C_0(\R^2;\R^2)$, $v\in L^2$, $\Div u=0$,
and $\mu=\curl u\in\M$. Then $u=K*\mu$.
\end{lemma}
\begin{proof}
The distribution $h=u-K*\mu$ has zero divergence and curl, so each component
is harmonic; Weyl's lemma makes it a smooth harmonic function. Its mean-value
estimate on $B_R(x)$ gives
\[
 |h(x)|\leq\frac1{|B_R|}\int_{B_R(x)}|U|
 +\frac{C\norm v2}{R}
 +\frac1{|B_R|}\int_{B_R(x)}|K*\mu|.
\]
The last term is at most $C\norm\mu\M/R$, because
$\sup_y\int_{B_R(x)}|K(z-y)|\,dz\leq CR$. The first tends to zero since
$U\in C_0$: separate a fixed large compact set from its uniformly small
complement. Letting $R\to\infty$ proves $h=0$.
\end{proof}

\begin{theorem}[Global existence]\label{thm:global}
For every datum in the exact target stated at the start of this document,
there is a global velocity with all of its stated properties.
\end{theorem}
\begin{proof}
Proposition~\ref{prop:iteration} gives the global velocity, the exact weak
initial trace and formulation, and every-time finite measure curl with
locally uniform variation. At each time take $U=V_0$ and
$v=(V(t)-V_0)+W(t)+\sum_qw_q(t)\in L^2$ in
Lemma~\ref{lem:canonical}; at zero use the prescribed datum directly.
This proves the remaining canonical Biot--Savart identity at every time.
No finite-horizon solutions have to be reconciled or restarted.
\end{proof}

\end{document}
